% ------------------------------------------------
% METHOD
% ------------------------------------------------

\section{Method}
\label{sec:methodology}

\subsection{Human reference data}
\label{subsec:humandata}

The human baseline is Study~2 of \textcite{atari2023weird}, a cross-cultural
validation of the second Moral Foundations Questionnaire (MFQ-2) administered in
19 countries, whose raw data and official questionnaire translations the authors
make publicly available. From that sample we use the five countries whose
dominant language is Spanish: Argentina, Chile, Colombia, Mexico and Peru
($n = 1{,}026$). They were selected because they differ from one another
culturally while sharing a language, which allows the prompt language and the
country cue to be varied independently. The file as distributed already contains
only respondents who passed all three attention checks, so applying that
criterion removes nobody; all item responses are complete and lie on the intended
1--5 scale.

\subsection{Instrument and scoring}
\label{subsec:instrument}

The MFQ-2 consists of 36 items measuring six moral foundations with six items
each: Care, Equality, Proportionality, Loyalty, Authority and Purity. Respondents
rate how well each statement describes them from \emph{Does not describe me at
all}~(1) to \emph{Describes me extremely well}~(5). A foundation score is the
mean of its six items and so also lies between 1 and 5.

The public data file contains 50 item columns rather than 36, because Study~2
administered a larger pool from which \textcite{atari2023weird} discarded 14
items for cross-loading. Matching the published items to columns by wording alone
is not sufficient, since three Loyalty items were reworded between the pool and
the final instrument. We therefore took the scoring key from the authors' own
Study~2 analysis script; the resulting mapping is given in
Appendix~\ref{app:mapping}.

\subsection{Design}
\label{subsec:design}

The experiment crosses \textbf{prompt language} (English or Spanish) with
\textbf{persona} (none, or one of the five countries), giving $2 \times 6 = 12$
cells per model. Following \textcite{munker2025cultural}, each cell is
administered 50 times independently, yielding 600 administrations per model and
1{,}800 across the three models. One \textbf{administration} is one complete
questionnaire delivered in a single model call: the model receives all 36 items
and returns all 36 answers, so the six foundation scores of an administration
come from one response. An administration is thus the synthetic counterpart of a
single human respondent.

\subsection{Prompts}
\label{subsec:prompts}

The instruction sentence, the 36 items and the five scale labels are taken
verbatim from the official English and Spanish MFQ-2 translations distributed
with \textcite{atari2023weird}, and were extracted programmatically from those
files rather than retyped. Using the official Spanish translation rather than
translating the items ourselves follows \textcite{aksoy2024whose} and ensures the
two language conditions administer the same instrument.

Each prompt has a system message and a user message. The system message carries
the persona and names only the country, for example \emph{``You are a person from
Mexico.''} In the no-persona baseline it is present but contentless and
structurally parallel (\emph{``You are a person answering a survey.''}), so that
what varies across persona conditions is the country cue rather than the presence
of a system message. The user message contains the official instruction, the five
labelled scale points, the 36 items in published order, and an instruction to
answer with a single JSON object mapping each item number to an integer from 1 to
5. Keeping the persona to a bare geographical cue follows
\textcite{munker2025cultural}, who describe such a prompt but do not print its
wording; we therefore implement the description and claim comparability of design
rather than of wording. Spanish strings not taken from the official translation
were back-translated with DeepL as an independent check and use the formal
\emph{usted} throughout, matching the register of the official instrument
(Appendix~\ref{app:translation}).

\subsection{Models and generation settings}
\label{subsec:models}

We use the three open-weight model families examined by
\textcite{munker2025cultural}, at the smaller size in each family, run locally
through Ollama~0.30.6 (Table~\ref{tab:models}). Running local open-weight models
keeps the families comparable with that study and fixes the system under test,
since the weights are identified by a content digest and cannot change beneath
the experiment as a hosted model can.

\begin{table}[H]
\centering
\caption{Models under test. Digests identify the exact weights; all three are
distributed by default at 4-bit quantisation.}
\label{tab:models}
\begin{tabular}{@{}llll@{}}
\toprule
\textbf{Model tag} & \textbf{Parameters} & \textbf{Quantisation} & \textbf{Digest (first 12)} \\
\midrule
\texttt{qwen2.5:7b}  & 7.6B & Q4\_K\_M & \texttt{845dbda0ea48} \\
\texttt{llama3.1:8b} & 8.0B & Q4\_K\_M & \texttt{46e0c10c039e} \\
\texttt{mistral:7b}  & 7.2B & Q4\_K\_M & \texttt{6577803aa9a0} \\
\bottomrule
\end{tabular}
\end{table}

Temperature was set explicitly to 0.8 rather than inherited from the runtime
default, the context window to 4{,}096 tokens and the generation limit to 1{,}024
tokens; \texttt{top\_p} was left unset. Each administration uses its own
deterministic seed derived from the model, language, persona and sample index, so
the full set of responses can be regenerated exactly.

\subsection{Response validation}
\label{subsec:validation}

A response is valid only if it parses as a JSON object containing exactly the
keys 1 to 36, each mapped to an integer between 1 and 5. Responses wrapped in
code fences or preceded by prose are recovered before validation; truncated,
incomplete or out-of-range responses are rejected and retried up to three times
with a fresh seed. Every raw response is stored whether or not it parsed, so no
response is discarded by a parsing decision. Parse and retry rates are treated as
results rather than diagnostics and are reported in Section~\ref{sec:results}.

\subsection{Metrics}
\label{subsec:metrics}

Let $m_c$ be a model's mean score on a given foundation under the persona for
country $c$, and $h_c$ the corresponding human mean, over the five countries $C$.

\textbf{Distance to humans} measures how far a model's answers sit from the
people it stands in for:
\begin{equation}
D = \frac{1}{|C|}\sum_{c \in C} \left| m_c - h_c \right|.
\end{equation}

\textbf{Between-country differentiation} asks whether the persona separates the
countries at all. We report the standard deviation of the five country means for
model and humans alike, and the effect size of a one-way ANOVA of persona on the
model's administrations, $\eta^2 = F(k-1)/\left(F(k-1)+(N-k)\right)$ with $k=5$.
The human country effect is reported beside it, so the model's differentiation is
judged against how much real respondents differ rather than against zero.

\textbf{Ordering accuracy} asks whether that separation runs in the right
direction, as Spearman's $\rho$ between the model's and the humans' ordering of
the five countries. Five countries admit only $5!=120$ orderings, which bounds
what a single foundation-level test can show (Section~\ref{sec:limitations}). We
therefore treat per-foundation $\rho$ as a descriptive effect size and test
ordering once per model and language using the mean $\rho$ across foundations,
against a null formed by permuting country labels independently within each
foundation (20{,}000 permutations).

\textbf{Within-country flattening} is the ratio of the model's standard deviation
across its 50 administrations for a country to the human within-country standard
deviation; a ratio below 1 indicates less variation than among real respondents.

\textbf{Language effect} compares, for the no-persona cells, each model's mean
under the two prompt languages and its absolute distance to the pooled human mean,
asking whether prompting in Spanish moves a model towards or away from
Spanish-speaking respondents.

Within each family of tests, $p$-values are adjusted using the Holm--Bonferroni
procedure, and exact permutation $p$-values are used for Spearman's $\rho$, whose
asymptotic approximation is not valid at $n=5$.
